\section{Comparative Evaluation}
\label{sec:eval}
\subsection{Setup and Metrics}
\label{sec:setup}
The November 16 evaluation contains 1,315 LDraw exports across two cohorts spanning 21 categories and five description levels. Group B is the main comparison: 265 shared prompts, 53 per level, with one delivered output per system per prompt. The BrickForma inputs are raw outputs generated with its Q8\_0 weights. Group A contains 260 matched prompts for BrickGPT and BrickNet (Appendix~\ref{app:eval}). We score all supplied exports in both cohorts.

We use uniform yellow material and render each export at $512{\times}512$ from eight orthographic azimuths at $45^\circ$ intervals and $25^\circ$ elevation. For each text--image metric we take the maximum over views, then average over prompts. VQAScore~\cite{vqascore} uses the first-token probability of \textquotedblleft Yes\textquotedblright\ from Qwen2.5-VL-7B-Instruct~\cite{qwen25vl}. PE~\cite{pe} uses cosine similarity of normalized embeddings, and SigLIP2~\cite{siglip2} uses the sigmoid of the image--text logit. We compute $P_{\mathrm{inv}}$ as the mean fraction of invalid placements per model in export order, using BrickNet's inset collision kernel~\cite{bricknet}.

\subsection{Performance across Description Levels}
Table~\ref{tab:levels} addresses Q1 by separating Group B into L1--L5. BrickForma leads PE and SigLIP2 at all five levels and VQAScore at L2--L5. At L1 (13.4 words on average), BrickGPT has a slightly higher VQAScore (0.1911 versus 0.1841). BrickGPT has the lowest mean $P_{\mathrm{inv}}$ at each level. Figure~\ref{fig:gallery} illustrates individual outcomes from the same comparison.
\begin{table*}[t]
\centering\small
\setlength{\tabcolsep}{3pt}
\begin{tabular}{lrrlrrrr}
\toprule
Level & $n$ & Words & System & $P_{\mathrm{inv}}\downarrow$ & VQAScore $\uparrow$ & PE $\uparrow$ & SigLIP2 $\uparrow$ \\
\midrule
L1 & 53 & 13.4 & BrickForma & 0.41\% & 0.1841 & \textbf{0.1409} & \textbf{0.0254} \\
 &  &  & BrickGPT & \textbf{0.00\%} & \textbf{0.1911} & 0.1323 & 0.0132 \\
 &  &  & BrickNet & 15.14\% & 0.0445 & 0.1032 & 0.0002 \\
\midrule
L2 & 53 & 26.4 & BrickForma & 0.78\% & \textbf{0.1651} & \textbf{0.1453} & \textbf{0.0241} \\
 &  &  & BrickGPT & \textbf{0.00\%} & 0.1354 & 0.1391 & 0.0206 \\
 &  &  & BrickNet & 17.74\% & 0.0579 & 0.1037 & 0.0013 \\
\midrule
L3 & 53 & 37.3 & BrickForma & 0.73\% & \textbf{0.1467} & \textbf{0.1479} & \textbf{0.0717} \\
 &  &  & BrickGPT & \textbf{0.00\%} & 0.1061 & 0.1396 & 0.0199 \\
 &  &  & BrickNet & 23.30\% & 0.0211 & 0.1075 & 0.0012 \\
\midrule
L4 & 53 & 52.3 & BrickForma & 1.38\% & \textbf{0.0239} & \textbf{0.1792} & \textbf{0.0914} \\
 &  &  & BrickGPT & \textbf{0.00\%} & 0.0198 & 0.1609 & 0.0487 \\
 &  &  & BrickNet & 18.68\% & 0.0023 & 0.1344 & 0.0016 \\
\midrule
L5 & 53 & 78.6 & BrickForma & 0.75\% & \textbf{0.0198} & \textbf{0.1727} & \textbf{0.0698} \\
 &  &  & BrickGPT & \textbf{0.00\%} & 0.0107 & 0.1575 & 0.0286 \\
 &  &  & BrickNet & 14.56\% & 0.0035 & 0.1379 & 0.0035 \\
\bottomrule
\end{tabular}
\caption{Group B means by description level. Each level contains 53 prompts per system. Words is the mean English word count. Bold marks the best mean.}
\label{tab:levels}
\end{table*}


\subsection{Aggregate Results}
\begin{table}[t]
\centering\small
\setlength{\tabcolsep}{3pt}
\begin{tabular}{lrrrr}
\toprule
System & $P_{\mathrm{inv}}\downarrow$ & VQAScore $\uparrow$ & PE $\uparrow$ & SigLIP2 $\uparrow$ \\
\midrule
BrickForma & 0.81\% & \textbf{0.1079} & \textbf{0.1572} & \textbf{0.0565} \\
BrickGPT & \textbf{0.00\%} & 0.0926 & 0.1459 & 0.0262 \\
BrickNet & 17.88\% & 0.0259 & 0.1173 & 0.0016 \\
\bottomrule
\end{tabular}
\caption{Group B means on 265 matched prompts. BrickForma uses delivered raw outputs. Bold marks the best mean.}
\label{tab:main}
\end{table}

BrickForma attains the highest mean VQAScore (0.1079), PE (0.1572), and SigLIP2 (0.0565) on Group B (Table~\ref{tab:main}). It exceeds BrickGPT on 188, 166, and 168 of 265 prompts for the three metrics, respectively. All six pairwise comparisons pass Holm correction under both paired tests (Appendix~\ref{app:eval}). The raw BrickForma exports have mean $P_{\mathrm{inv}}$ of 0.81\%, compared with 0.00\% for BrickGPT and 17.88\% for BrickNet.

\subsection{Geometry Before and After Reconstruction}
\label{sec:diagnostics}
We assess reconstruction on 110 raw/retiled pairs from the earlier 114-prompt comparison. Adjacent-layer footprint overlap defines stud contacts. After mapping exports to a grid with 20 LDU horizontal pitch and 24 LDU layer height and translating occupied cells to a common origin, all 110 pairs have equal voxel unions. The number with a single stud component rises from 18 to 110 (Table~\ref{tab:export-audit}).
\begin{table}[t]
\centering\small
\begin{tabular}{@{}lrr@{}}
\toprule
Export measure & Raw & Retiled \\
\midrule
Single stud component & 18/110 & 110/110 \\
Mean stud components & 4.29 & 1.00 \\
No body-cell overlap & 100/110 & 110/110 \\
\midrule
\multicolumn{3}{@{}l@{}}{Equal voxel unions after translation: 110/110} \\
\bottomrule
\end{tabular}
\caption{Earlier export audit of 110 matched raw/retiled pairs. Union equality uses translation normalization.}
\label{tab:export-audit}
\end{table}

Across these pairs $P_{\mathrm{inv}}$ falls from 0.1897\% to 0.0000\%, while mean part count rises from 178.7 to 421.2. In the separate 525-prompt diagnostic cohort, 46 raw layouts (8.8\%) have a single stud component, and all 525 final layouts pass both geometry checks after the reported recovery procedure (Appendix~\ref{app:eval}).

\section{Limitations and Discussion}
The transfer described in Q2 combines program design, geometric checks, and supervised fine-tuning. Isolating the contribution of large-model program authoring would require controls with alternative generators and matched training budgets.

The delivered systems use different inference settings and part libraries. Description levels vary in both content and length, and long descriptions may be truncated (Appendix~\ref{app:eval}). Alignment scores describe rendered appearance, while $P_{\mathrm{inv}}$ measures collisions under the stated file order. Physical stability and assembly-order feasibility require separate validation.

\section{Conclusion}
BrickForma uses executable programs to synthesize voxel shapes and converts them into checked training layouts for Qwen3-1.7B. On the 265 matched prompts, its raw exports lead all three alignment metrics in aggregate. Earlier reconstruction audits show that exact-cover repair can preserve the predicted volume while supplying connections that the model frequently misses.
